% ============================================================================
% Relazione del progetto di Programmazione a Oggetti (a.a. 2025/26)
% "Attivita": gestione di attivita' personali con C++ e Qt
%
% Prima di consegnare:
%  1. compilare i campi \studente e \matricola;
%  2. inserire il diagramma UML in figura/uml_attivita.png;
%  3. aggiornare la tabella delle ore (sezione 7) con le ore effettive;
%  4. compilare: pdflatex relazione.tex
% ============================================================================
\documentclass[10pt,a4paper]{article}

\usepackage[T1]{fontenc}
\usepackage[utf8]{inputenc}
\usepackage[a4paper,margin=2.2cm]{geometry}
\usepackage{graphicx}
\usepackage{booktabs}
\usepackage{array}
\usepackage{hyperref}
\usepackage{enumitem}
\usepackage{xcolor}

% Caption in italiano (la lingua babel "italian" non e' disponibile ovunque)
\renewcommand{\figurename}{Figura}
\renewcommand{\tablename}{Tabella}

\hypersetup{colorlinks=true, linkcolor=black, urlcolor=blue}

% ---- dati dello studente (da compilare) ------------------------------------
\newcommand{\studente}{Nome Cognome}
\newcommand{\matricola}{MATRICOLA}

\setlength{\parskip}{4pt}

\begin{document}

\begin{center}
{\LARGE\bfseries Attivita: gestione di attivita' personali}\\[4pt]
{\large Applicazione desktop in C++ con interfaccia grafica Qt}\\[8pt]
{\normalsize Progetto di Programmazione a Oggetti -- a.a.\ 2025/26}\\[4pt]
{\normalsize \studente\ --- matricola \matricola\ (progetto individuale)}
\end{center}

\vspace{6pt}
\noindent\hrulefill
\vspace{6pt}

\section{Introduzione}

L'obiettivo del progetto e' sviluppare un'applicazione per la gestione di
attivita' personali -- impegni, eventi, scadenze e promemoria -- interamente
scritta in C++ con interfaccia grafica realizzata con il framework Qt. Gli
utenti possono creare, modificare, eliminare, cercare e visualizzare i diversi
tipi di attivita', ciascuno dei quali possiede attributi specifici: ogni tipo e'
rappresentato da una classe derivata da una classe base comune che descrive gli
attributi e i comportamenti condivisi.

L'architettura del progetto separa nettamente il modello logico dall'interfaccia
grafica:

\begin{itemize}[nosep]
  \item \textbf{modello logico} (\texttt{model/}): una libreria C++20 pura
    (namespace \texttt{events}), riutilizzabile senza dipendenze dall'interfaccia,
    che implementa la gerarchia delle attivita', la generazione delle occorrenze
    temporali e la collezione \texttt{Calendar};
  \item \textbf{persistenza} (\texttt{model/persistence/}): un modulo separato
    (Qt Core, niente Widgets) che serializza e ricarica le attivita' su file
    locale in formato JSON tramite il pattern Visitor;
  \item \textbf{interfaccia grafica} (\texttt{app/}): un'applicazione Qt a
    finestra singola (pattern MVC) che usa il modello e la persistenza.
\end{itemize}

Il programma compila senza errori nel container di valutazione del corso
(Qt 6.4.2) sia con \texttt{cmake} sia con \texttt{qmake} (\texttt{Attivita.pro}),
ed esegue in modalita' standalone: i dati sono salvati su file JSON locale,
scelto dall'utente tramite finestre di dialogo (mai percorsi cablati).

\section{Classi principali del modello logico}

La Fig.~\ref{fig:uml} mostra il diagramma UML delle classi principali del
modello. La radice della gerarchia e' la classe astratta \texttt{Activity},
che espone l'interfaccia polimorfa comune: il titolo, il punto temporale di
riferimento, il metodo \texttt{occurrencesIn}, la clonazione (\texttt{clone}),
il doppio dispatch (\texttt{accept}) e una descrizione testuale
(\texttt{describe}).

\begin{figure}[h]
\centering
\IfFileExists{figura/uml_attivita.png}{%
  \includegraphics[width=0.92\textwidth]{figura/uml_attivita.png}%
}{%
  \fbox{\parbox[c][4cm][c]{0.9\textwidth}{\centering Diagramma UML da inserire}}%
}
\caption{Diagramma UML delle classi principali del modello logico}
\label{fig:uml}
\end{figure}

\subsection{La gerarchia delle attivita'}

Il progetto e' individuale, quindi e' richiesta una gerarchia con almeno tre
classi concrete: ne sono implementate \textbf{quattro}, ciascuna con attributi
e metodi genuinamente diversi (Tab.~\ref{tab:tipi}).

\begin{table}[h]
\centering
\begin{tabular}{@{}l l l@{}}
\toprule
\textbf{Classe} & \textbf{Attributi specifici} & \textbf{Metodi specifici} \\
\midrule
\texttt{Event} & inizio, durata & \texttt{getEnd}, \texttt{isIn}, \texttt{overlapsWith} \\
\texttt{RecurrentEvent} & generatore, template, eccezioni & \texttt{addException}, \texttt{truncateBefore} \\
\texttt{Deadline} & scadenza, priorita', evasa & \texttt{isOverdue}, \texttt{timeRemaining}, \texttt{setDone} \\
\texttt{Reminder} & attivazione, messaggio, ripetizione & \texttt{snooze}, \texttt{isRepeating} \\
\bottomrule
\end{tabular}
\caption{I quattro tipi concreti di attivita' e i loro attributi/metodi distintivi}
\label{tab:tipi}
\end{table}

\textbf{Event} e' l'attivita' con un unico intervallo temporale: titolo, inizio
e durata (la durata negativa e' rifiutata sia nel costruttore sia in
\texttt{setDuration}). \textbf{RecurrentEvent} e' l'attivita' ricorrente: una
regola di ricorrenza (un \texttt{DateGenerator}), un evento \emph{template} da
cui derivano le occorrenze e un insieme di eccezioni (semantica EXDATE
dello standard iCalendar): le eccezioni \emph{rimuovono} occorrenze che il
generatore produrrebbe, e vivono sull'evento, non sul generatore.
\texttt{truncateBefore} termina la serie a partire da una data. \textbf{Deadline}
rappresenta una scadenza puntuale con priorita' (bassa/media/alta), stato di
completamento e i metodi \texttt{isOverdue}/\texttt{timeRemaining}.
\textbf{Reminder} e' un promemoria con messaggio, ripetizione opzionale e
\texttt{snooze} (posticipo dell'attivazione).

\subsection{Generatori di date (Strategy)}

La classe astratta \texttt{DateGenerator} produce i punti temporali delle
occorrenze e viene usata da \texttt{RecurrentEvent} per delega (pattern
\textbf{Strategy}): la regola di ricorrenza e' un oggetto separato, sostituibile
e mutabile a runtime senza toccare l'evento. Le implementazioni sono
\texttt{FixedIntervalGenerator} (ogni $N$ secondi, con fine opzionale),
\texttt{YearlyGenerator} (una volta all'anno, con gestione degli anni bisestili:
il 29 febbraio viene spostato al 28) e \texttt{NullGenerator} (null object).
Gli intervalli di interrogazione sono \emph{inclusivi} su entrambi gli estremi.

\subsection{Occurrence e Calendar}

\texttt{Occurrence} e' un value type leggero
(\texttt{\{source, start, duration\}}) che rappresenta una singola occorrenza
di un'attivita' sulla timeline; \texttt{Calendar} e' la collezione polimorfa
che possiede le attivita' (\texttt{unique\_ptr}), ne aggrega le occorrenze
ordinate (\texttt{occurrencesIn}) e offre la ricerca testuale per titolo
(\texttt{search}, case-insensitive). La creazione degli oggetti passa
esclusivamente da \texttt{ActivityFactory} (pattern \textbf{Factory}).

\section{Utilizzo non banale del polimorfismo}

Il polimorfismo e' utilizzato in modo non banale in quattro punti, riassunti
nella Tab.~\ref{tab:poly}; nessuna classe del modello espone un metodo
\texttt{getType} che restituisca una stringa per il controllo di flusso.

\begin{table}[h]
\centering
\begin{tabular}{@{}p{3.2cm} p{7.2cm} p{4.2cm}@{}}
\toprule
\textbf{Metodo polimorfo} & \textbf{Comportamento per tipo dinamico} & \textbf{Valore aggiunto} \\
\midrule
\texttt{occurrencesIn(from, to)} &
\texttt{Event} $\rightarrow$ 0/1 occorrenza; \texttt{RecurrentEvent} $\rightarrow$
espansione della ricorrenza meno le eccezioni; \texttt{Deadline}/\texttt{Reminder}
$\rightarrow$ istante/i puntuale/i &
La timeline (vista settimanale, aggiornamento automatico) interroga tutte le
attivita' allo stesso modo; l'espansione temporale e' incapsulata in ogni classe \\
\addlinespace
\texttt{accept(visitor)} (doppio dispatch) &
\texttt{ActivityVisitor::visit(const Event\&)}, \ldots/\texttt{Deadline}/
\texttt{Reminder} &
Persistenza JSON e pannelli GUI (dettaglio, etichette) aggiunti senza
modificare le classi del dominio (Single Responsibility) \\
\addlinespace
\texttt{clone()} (Prototype) &
Copia profonda del tipo concreto, con ritorno covariante &
\texttt{RecurrentEvent} genera occorrenze come \emph{copie indipendenti} del
template, modificabili senza effetti collaterali \\
\addlinespace
\texttt{generateDates} (Strategy) &
Ogni \texttt{DateGenerator} produce le date a modo suo (intervallo fisso,
annuale, mai) &
La regola di ricorrenza si sostituisce e si modifica a runtime
(\texttt{truncateBefore} $\rightarrow$ \texttt{setEnd}) \\
\bottomrule
\end{tabular}
\caption{Metodi polimorfi non banali e loro valore aggiunto}
\label{tab:poly}
\end{table}

Il caso piu' rilevante e' \texttt{occurrencesIn}: il metodo si comporta in modo
profondamente diverso in base al tipo dinamico dell'oggetto invocante ed e'
utilizzato da tutta l'applicazione (vista settimanale, persistenza, test).
Il doppio dispatch (\texttt{accept}) e' la base del pattern Visitor utilizzato
per la persistenza (sezione 4) e per la costruzione dei pannelli di dettaglio
nella GUI (sezione 6), consentendo di aggiungere operazioni esterne alla
gerarchia senza toccare le classi del modello e senza \texttt{dynamic\_cast}.

\section{Persistenza dei dati}

Conformemente ai vincoli del corso, i dati delle attivita' sono salvati su un
\textbf{file locale in formato JSON} (formato strutturato), scelto dall'utente
in ogni momento tramite finestre di dialogo grafiche (\texttt{QFileDialog}):
nessun percorso e' cablato nel codice sorgente.

La serializzazione e' realizzata con il \textbf{pattern Visitor}: le classi del
dominio espongono \texttt{accept(ActivityVisitor\&)} e la persistenza e'
implementata in due visitor esterni, \texttt{JsonActivityVisitor} (per le
attivita') e \texttt{JsonGeneratorVisitor} (per la regola di ricorrenza,
attraverso una seconda gerarchia di visitor). In questo modo le classi del
modello non contengono alcuna logica JSON e nuovi formati di export potranno
essere aggiunti in futuro implementando soltanto nuovi visitor.

La deserializzazione e' una \emph{factory} con discriminatore: il campo
\texttt{"type"} dell'oggetto JSON (\texttt{event}, \texttt{recurrent},
\texttt{deadline}, \texttt{reminder}) determina la classe concreta da
ricostruire. Ogni campo e' validato rigorosamente: date ISO-8601
(\texttt{YYYY-MM-DDTHH:MM:SS}, formato UTC, lunghezza esatta e date possibili),
durate non negative, intervalli di ricorrenza positivi, priorita' note; in
presenza di dati non validi il caricamento fallisce con un messaggio d'errore
e il contenuto del calendario non viene alterato. Il formato del file e':

\begin{verbatim}
{ "version": 1, "activities": [ { "type": "recurrent",
    "template": { "title": "Lezione", "start": "2026-09-07T14:00:00",
                  "duration_seconds": 7200 },
    "generator": { "type": "fixed", "start": "2026-09-07T14:00:00",
                   "interval_seconds": 604800, "end": "2026-12-21T00:00:00" },
    "exceptions": [ "2026-11-16T14:00:00" ] }, ... ] }
\end{verbatim}

Il modulo di persistenza dipende solo da \textbf{Qt Core} (per la gestione
JSON) e non dall'interfaccia grafica: il modello logico rimane riutilizzabile
senza dipendenze dalla GUI, come richiesto dai vincoli del corso. Il file
d'esempio \texttt{examples/attivita\_esempio.json}, fornito nella consegna,
e' stato generato e validato dalla libreria stessa (round-trip completo).

\section{Interfaccia grafica}

L'applicazione Qt (\texttt{app/}) adotta il pattern \textbf{MVC}: il
\texttt{CalendarController} possiede il \texttt{Calendar} logico ed espone le
operazioni (creazione, modifica, eliminazione, ricerca, salvataggio,
caricamento); un segnale \texttt{activitiesChanged} notifica alle viste ogni
modifica, che si aggiornano automaticamente. La finestra principale e'
organizzata in tre pagine navigabili tramite \texttt{QStackedWidget}:

\begin{enumerate}[nosep]
  \item \textbf{settimana}: griglia settimanale stile Google Calendar disegnata
    con \texttt{QPainter} (giorni, ore, occorrenze come blocchi colorati, tooltip,
    menu contestuale, navigazione settimane); le occorrenze che si sovrappongono
    temporalmente nello stesso giorno vengono \textbf{affiancate in colonne}
    (coesistono nella stessa casella); \textbf{trascinamento}: tenendo premuta
    un'occorrenza la si sposta (drag\&drop) su una nuova data/ora, aggiornando i
    dati dell'attivita' (per i ricorrenti l'intera serie, eccezioni e fine
    comprese); doppio clic su una cella vuota avvia la creazione precompilata;
  \item \textbf{elenco}: tabella delle attivita' con ricerca live per titolo;
  \item \textbf{dettaglio}: campi specifici per tipo, costruiti con un Visitor
    (la GUI non discrimina mai tramite stringhe di tipo).
\end{enumerate}

La creazione e la modifica avvengono in un \textbf{pannello ridotto mostrato
dentro la finestra principale}: piu' piccolo della finestra, e' un widget
figlio della MainWindow e pertanto \emph{non puo' uscire dai suoi bordi}: la
sua dimensione viene limitata a quella della finestra (\texttt{boundedTo}), e
viene ricentrato automaticamente anche quando la finestra viene ridimensionata.
Al suo interno i campi sono organizzati in righe con l'etichetta a sinistra e la
box di input affiancata (es. \emph{Titolo} con il relativo campo di testo,
\emph{Data} con la casella data/ora). L'utente sceglie il tipo da un menu a
tendina (\texttt{Evento}, \texttt{Evento ricorrente}, \texttt{Scadenza},
\texttt{Promemoria}) e si compila il pannello con i campi specifici di quel
tipo: titolo e data/ora in tutti i tipi (conservati anche cambiando tipo);
durata per gli eventi; intervallo di ripetizione e fine opzionale (tramite
checkbox) per i ricorrenti; priorita' (bassa/media/alta) e stato "evasa" per le
scadenze; messaggio di dettaglio e ricorrenza per i promemoria. Il doppio clic
su una cella vuota della settimana precompila data e ora dell'inizio nel modulo;
il doppio clic su un'occorrenza esistente apre invece la \textbf{modifica
dell'attivita'} conservando il tipo originale (per i ricorrenti l'intera serie,
con regola e data di fine modificabili); il menu contestuale offre anche la
modifica di una singola istanza. Il pannello si chiude con i pulsanti
Salva/Annulla.

La toolbar permette di creare/modificare/eliminare, di passare tra le viste e
di \textbf{salvare e caricare} il file JSON in qualunque momento tramite
\texttt{QFileDialog}. La modifica di una singola occorrenza di un'attivita'
ricorrente aggiunge un'eccezione (EXDATE) e crea un evento singolo sostitutivo;
l'eliminazione di un'istanza offre le scelte ``solo questa'' (eccezione) o
``questa e le successive'' (troncamento della ricorrenza).

\section{Funzionalita' aggiuntive}

\begin{itemize}[nosep]
  \item gestione di \textbf{quattro tipi} di attivita' (il minimo richiesto e'
    tre), con attributi e metodi specifici;
  \item \textbf{ricorrenze} con eccezioni (EXDATE), troncamento e ricorrenza
    annuale (compleanni) con gestione degli anni bisestili;
  \item \textbf{vista settimanale personalizzata} disegnata con \texttt{QPainter}
    (scaling proporzionale, evidenza del giorno corrente, tooltip, menu
    contestuale, affiancamento degli eventi sovrapposti, \textbf{drag\&drop}
    per spostare rapidamente le attivita');
  \item \textbf{ricerca live} per titolo (case-insensitive) nell'elenco;
  \item \textbf{modifica di una singola istanza} di una ricorrenza senza
    alterare le altre;
  \item \textbf{persistenza robusta}: caricamento con validazione completa del
    formato e messaggi d'errore; file d'esempio incluso nella consegna;
  \item suite di \textbf{test automatici} (framework Catch2): oltre venti casi
    su modello e persistenza e sul controller dell'applicazione (round-trip
    JSON per ogni tipo, eccezioni, troncamenti, validazioni); i test sono
    eseguiti nella pipeline Docker del repository.
\end{itemize}

\section{Rendicontazione delle ore}

Il progetto e' svolto individualmente. Le ore effettive sono riportate nella
colonna dedicata.

\begin{table}[h]
\centering
\begin{tabular}{@{}l c c@{}}
\toprule
\textbf{Attivita'} & \textbf{Previste} & \textbf{Effettive} \\
\midrule
Analisi della specifica e progettazione & 6h & ~h \\
Modello logico (gerarchia \texttt{Activity}, generatore di date) & 10h & ~h \\
Persistenza JSON (Visitor) & 6h & ~h \\
Interfaccia grafica Qt (MVC, viste, form, ricerca) & 12h & ~h \\
Verifica nel container del corso, build e test & 3h & ~h \\
Relazione e disegni UML & 3h & ~h \\
\midrule
\textbf{Totale} & \textbf{40h} & \textbf{~h} \\
\bottomrule
\end{tabular}
\caption{Rendicontazione delle ore previste ed effettive}
\label{tab:ore}
\end{table}

\end{document}
